% Basic commands
\usepackage[margin=1.5in]{geometry}
\usepackage[spanish]{babel}
\usepackage{amsmath, amsfonts, mathtools, amsthm, amssymb}
\usepackage[usenames,dvipsnames]{xcolor}
\usepackage{float}
\usepackage{microtype}
\usepackage{enumitem}

% Tikz
\usepackage{tikz}
\usepackage{tikz-cd}
\usetikzlibrary{intersections, angles, quotes, calc, positioning, arrows.meta, shadows}
\usepackage{pgfplots}
\pgfplotsset{compat=1.13}

% For math environments
\usepackage{thmtools}
\usepackage[framemethod=TikZ]{mdframed}
\mdfsetup{skipabove=1em,skipbelow=0.6em}
\theoremstyle{definition}

%%%%%%%%%%%%%%%%%%%%%%%%%%%%%%%%%%%%%%%%%%%%%%%%%%%%%%%%%%%%
% STYLE 1: Midnight Atlas
% Dark academic blue with rounded cards and soft shadows.
%%%%%%%%%%%%%%%%%%%%%%%%%%%%%%%%%%%%%%%%%%%%%%%%%%%%%%%%%%%%
\definecolor{MidnightInk}{HTML}{10233F}
\definecolor{AtlasBlue}{HTML}{1E4E8C}
\definecolor{AtlasSoft}{HTML}{EEF5FF}
\definecolor{AtlasLine}{HTML}{527BB5}
\definecolor{AtlasGold}{HTML}{B88A2B}

\declaretheoremstyle[
    headfont=\bfseries\sffamily\color{MidnightInk},
    bodyfont=\normalfont,
    %spaceabove=0.8em,
    %spacebelow=0.8em,
    mdframed={
        nobreak,
        roundcorner=9pt,
        linewidth=1.1pt,
        linecolor=AtlasBlue!85!black,
        backgroundcolor=AtlasSoft,
        %innerleftmargin=11pt,
        %innerrightmargin=11pt,
        %innertopmargin=9pt,
        %innerbottommargin=10pt,
        shadow=true,
        shadowsize=4pt,
        shadowcolor=black!18,
    }
]{atlasbox}

\declaretheoremstyle[
    headfont=\bfseries\sffamily\color{AtlasGold!70!black},
    bodyfont=\normalfont,
    %spaceabove=0.7em,
    %spacebelow=0.7em,
    mdframed={
        nobreak,
        roundcorner=7pt,
        linewidth=0.9pt,
        linecolor=AtlasGold!70!black,
        backgroundcolor=AtlasGold!7,
        %innerleftmargin=10pt,
        %innerrightmargin=10pt,
        %innertopmargin=8pt,
        innerbottommargin=9pt,
        shadow=true,
        shadowsize=3pt,
        shadowcolor=black!12,
    }
]{atlasdef}

\declaretheoremstyle[
    headfont=\bfseries\sffamily\color{MidnightInk!80},
    bodyfont=\normalfont,
    mdframed={
        linewidth=0pt,
        backgroundcolor=black!3,
        roundcorner=6pt,
        innerleftmargin=10pt,
        innerrightmargin=10pt,
        innertopmargin=7pt,
        innerbottommargin=8pt,
    }
]{atlasserene}

\declaretheoremstyle[
    headfont=\bfseries\sffamily\color{AtlasBlue!80!black},
    bodyfont=\normalfont,
    mdframed={
        linewidth=1.1pt,
        linecolor=AtlasLine,
        rightline=false,
        topline=false,
        bottomline=false,
        innerrightmargin=8pt,
        innerleftmargin=9pt,
    }
]{atlasline}

\declaretheoremstyle[
    headfont=\bfseries\sffamily\color{BrickRed!60!black},
    bodyfont=\normalfont,
    mdframed={
        linewidth=1.1pt,
        linecolor=BrickRed!60!black,
        rightline=false,
        topline=false,
        bottomline=false,
        innerrightmargin=8pt,
        innerleftmargin=9pt,
    }
]{obsline}

\declaretheorem[style=atlasdef, name=Definición, numberwithin=chapter]{definition}
\declaretheorem[style=atlasserene, numbered=no, name=Ejemplo]{eg}
\declaretheorem[style=atlasbox, name=Proposición, numberwithin=chapter]{prop}
\declaretheorem[style=atlasbox, name=Teorema, numberwithin=chapter]{theorem}
\declaretheorem[style=atlasbox, name=Lema, numberwithin=chapter]{lema}
\declaretheorem[style=atlasbox, name=Corolario, numberwithin=chapter]{colorary}
\declaretheorem[style=atlasline, numbered=no, name=Notación]{notation}
\declaretheorem[style=obsline, numbered=no, name=Observación]{observation}

% Headers and footers
\usepackage{fancyhdr}
\pagestyle{fancy}
\fancyhf{}
\fancyhead[L]{\sffamily\nouppercase{Victoria Eugenia Torroja}}
\fancyhead[R]{\sffamily\nouppercase\rightmark}
\fancyfoot[C]{\sffamily\leftmark}
\fancyfoot[R]{\sffamily\thepage}
\renewcommand{\headrulewidth}{0.4pt}
\renewcommand{\footrulewidth}{0pt}

% New commands
\newcommand{\R}{\mathbb{R}}
\newcommand{\C}{\mathbb{C}}
\newcommand{\F}{\mathbb{F}}
\newcommand{\N}{\mathbb{N}}
\newcommand{\Q}{\mathbb{Q}}
\newcommand{\Z}{\mathbb{Z}}
\newcommand{\K}{\mathbb{K}}
\newcommand{\mcd}{\text{mcd}}
\newcommand{\mcm}{\text{mcm}}
\DeclareMathOperator{\Ker}{Ker}
\DeclareMathOperator{\Imagen}{Im}
\DeclareMathOperator{\Hom}{Hom}
\DeclareMathOperator{\Aut}{Aut}
\DeclareMathOperator{\ran}{ran}
\DeclareMathOperator{\GL}{GL}
\DeclareMathOperator{\SL}{SL}
\DeclareMathOperator{\SO}{SO}
\DeclareMathOperator{\ord}{ord}
\DeclareMathOperator{\ind}{ind}
\DeclareMathOperator{\sig}{sig}
\DeclareMathOperator{\dom}{dom}
\DeclareMathOperator{\End}{End}
\DeclareMathOperator{\Adj}{Adj}
\DeclareMathOperator{\grad}{grad}
\DeclareMathOperator{\traz}{traz}
\DeclareMathOperator{\arctanh}{arctanh}
\DeclareMathOperator{\arcsinh}{arcsinh}
\DeclareMathOperator{\arccosh}{arccosh}
\DeclareMathOperator{\rad}{rad}
\DeclareMathOperator{\ad}{ad}
\DeclareMathOperator{\Bil}{Bil}
\DeclareMathOperator{\sech}{sech}
\DeclareMathOperator{\Int}{Int}
\DeclareMathOperator{\Div}{div}
\DeclareMathOperator{\rot}{rot}
\DeclareMathOperator{\Rel}{Re}
\DeclareMathOperator{\Arg}{Arg}
\let\epsilon\varepsilon

